% ============================================================================
\chapter{Raspberry Pi deployment}
\label{ch:rpi}
% ============================================================================

This chapter describes how to deploy and run the \appname{} application on a
Raspberry~Pi with a 7\,'' touch screen (800\,$\times$\,480 resolution) as a
dedicated, touch-only operator terminal.

\section{Operating system}

Use \textbf{Raspberry Pi OS (64-bit)}. The Lite version is suitable --- and
preferred --- if you run the application in kiosk mode without a full desktop
environment: the application can render directly to the screen
(Section~\ref{sec:rendering}).

\section{System dependencies}

When building the application on the Pi, or running it with the default renderer,
install the following system libraries first:

\begin{lstlisting}[language=bash]
$ sudo apt-get update
$ sudo apt-get install -y build-essential cmake \
      libfreetype6-dev libfontconfig1-dev libssh-dev \
      libxkbcommon-dev libinput-dev libudev-dev libgbm-dev
\end{lstlisting}

\section{Building on the Pi and cross-compilation}

The straightforward approach is to build natively on the Pi~5 (a release build takes
a few minutes): copy or clone the repository onto the Pi and run
\code{cargo build --release}.

Building on the Pi can be slow on older models; you can instead cross-compile from a
faster machine using \code{cross-rs}:

\begin{lstlisting}[language=bash]
$ cargo install cross
$ cross build --target aarch64-unknown-linux-gnu --release
\end{lstlisting}

Then copy the resulting binary \emph{and} \code{config.toml} to the Pi.

\section{Performance and rendering}
\label{sec:rendering}

Slint supports multiple rendering backends. On a Raspberry~Pi, the
\textbf{\code{linuxkms}} backend is often the most performant, as it renders
directly to the screen without X11 or Wayland overhead.

To run with LinuxKMS:

\begin{enumerate}[itemsep=2pt]
  \item Ensure your user is in the \code{video} and \code{input} groups:
\begin{lstlisting}[language=bash]
$ sudo usermod -aG video,input,render $USER
\end{lstlisting}
  \item Run the application with the backend selected via the environment:
\begin{lstlisting}[language=bash]
$ SLINT_BACKEND=linuxkms ./deltax2
\end{lstlisting}
\end{enumerate}

\begin{note}
The \code{linuxkms} backend is an optional Slint feature and must be enabled at
compile time. If \code{SLINT\_BACKEND=linuxkms} reports an unknown backend, add the
feature to \code{Cargo.toml} and rebuild:
\begin{lstlisting}[language=TOML]
slint = { version = "1.14", features = ["backend-linuxkms"] }
\end{lstlisting}
When running under a desktop session (X11 or Wayland) instead, no environment
variable is needed; the default \code{winit} backend is used.
\end{note}

\section{Kiosk mode}

For daily operation the application should fill the screen without window
decorations. This is controlled by the configuration file
(Chapter~\ref{ch:configuration}):

\begin{lstlisting}[language=TOML]
[ui]
kiosk_mode = true
\end{lstlisting}

\section{Automatic start with systemd}
\label{sec:systemd}

To run the application as a standalone interface on boot, create the unit file
\code{deltax2.service} in \code{/etc/systemd/system/} (adjust user and paths):

\begin{lstlisting}
[Unit]
Description=DeltaX2 seeding robot control UI
After=multi-user.target

[Service]
User=pi
WorkingDirectory=/home/pi/deltax2
Environment=SLINT_BACKEND=linuxkms
ExecStart=/home/pi/deltax2/target/release/deltax2
Restart=on-failure
RestartSec=3

[Install]
WantedBy=multi-user.target
\end{lstlisting}

Then enable and start it:

\begin{lstlisting}[language=bash]
$ sudo systemctl daemon-reload
$ sudo systemctl enable --now deltax2.service
$ journalctl -u deltax2 -f        # follow the application log
\end{lstlisting}

All console messages of the application (connection results, executed serial
commands, port listings) appear in the systemd journal, which is the primary
diagnostic channel when working over SSH.

If you use a display manager and a desktop session instead of the bare-metal
\code{linuxkms} setup, an \code{.xsession} file or the desktop environment's
autostart mechanism can launch the application in the same way (without the
\code{SLINT\_BACKEND} variable).

\section{Touch screen calibration}
\label{sec:touchcal}

If the touch input is inverted or misaligned, configure \code{libinput} with a udev
rule that sets the \code{LIBINPUT\_CALIBRATION\_MATRIX}, e.g.\ in
\code{/etc/udev/rules.d/99-touch.rules}:

\begin{lstlisting}
ENV{ID_INPUT_TOUCHSCREEN}=="1", \
  ENV{LIBINPUT_CALIBRATION_MATRIX}="1 0 0 0 1 0"
\end{lstlisting}

The six values form the first two rows of a 3$\times$3 transformation matrix;
for example \code{-1 0 1 0 -1 1} rotates the input by 180°.

\section{Serial port permissions}

If the application cannot open the serial port, ensure the user has permission to
access serial devices:

\begin{lstlisting}[language=bash]
$ sudo usermod -aG dialout $USER
\end{lstlisting}

Log out and back in (or restart the systemd service) for the change to take effect.
